\chapter*{ГЛАВА 1}
\addcontentsline{toc}{chapter}{ГЛАВА 1}
\stepcounter{chapter}
\sloppy

\section{Анализ существующих курсов по машинному обучению}

Для формирования методологической и содержательной основы разрабатываемого курса был проведён анализ семи образовательных программ по машинному обучению. Поскольку область активно развивается, при отборе приоритет отдавался актуальным и востребованным программам с открытыми материалами, допускающими содержательный анализ системы контроля знаний.

Выборка охватывает три сегмента. Первый --- академические курсы ведущих российских университетов (НИУ ВШЭ, МГУ, МФТИ): они задают стандарт теоретической глубины и отражают опыт, с которым сталкивается студент профильной специальности. Второй --- открытые образовательные инициативы сообществ (ODS, girafe-ai): они представляют альтернативу университетскому образованию для широкой аудитории без формального отбора. Третий --- коммерческие онлайн-платформы (Karpov.Courses, Яндекс.Практикум): они отражают запросы рынка труда и ориентированы на практическую подготовку специалистов.

Для каждого курса рассматривались формат обучения, содержание программы, система контроля знаний и степень автоматизации процесса оценивания. Все анализируемые курсы ведутся на русском языке.

\subsection{Машинное обучение на ФКН НИУ ВШЭ}

Курс Е. А. Соколова~\cite{sokolov_hse} является одним из наиболее известных академических курсов по машинному обучению. Обучение проходит в очном формате для бакалавров факультета компьютерных наук НИУ ВШЭ, при этом все материалы, включая конспекты лекций и видеозаписи 2021 года, опубликованы в открытом доступе. Программа охватывает классические разделы машинного обучения.

Система контроля включает теоретические и практические домашние задания, контрольную работу и экзамен. Теоретические домашние задания и контрольные работы включают открытые вопросы с объяснением, а практические подразумевают работу в Jupyter Notebook. Домашние задания принимаются через платформу Anytask - она может только собирать ответ студентов, при этом не поддерживает автоматизацию проверки.

Практические задания в Jupyter Notebook частично автоматизированы юнит-тестами, однако все практические и теоретические работы проверяют ассистенты. Ручная проверка таких заданий фиксируется как ключевое ограничение: при росте потока студентов она масштабируется только увеличением количества ассистентов.

\subsection{Машинное обучение на ВМК МГУ}

Курс машинного обучения для студентов третьего курса факультета вычислительной математики и кибернетики МГУ~\cite{msu_ml} по содержанию близок к программе НИУ ВШЭ. Обучение организовано в очном формате и сопровождается практическими и теоретическими видами контроля. В практические входит работа в Jupyter Notebook и сореванования на платформе Kaggle, а в теоретические - коллоквиумы, тесты на лекциях, открытые вопросы в лабораторных работах и экзамен.

Проверка практических заданий автоматизирована при помощи юнит-тестов, при этом некоторые задания требуют кросс-проверки между студентами в количестве трех человек ~\cite{msu_ml_grading}.

Использование автотестов демонстрирует реализуемость автоматической проверки кода
прямо внутри Jupyter Notebook. Кросс-проверка открытых ответов тремя студентами, напротив,
отмечается как неудовлетворительный подход: она не обеспечивает
согласованности оценок и перекладывает нагрузку на самих студентов.

\subsection{Deep Learning School МФТИ}

Школа глубокого обучения (DLS) при МФТИ~\cite{dls_mipt} является бесплатным онлайн-курсом, реализованным на платформе Stepik. Программа разделена на четыре части. В первой части обсуждаются основы машинного обучения и нейронных сетей, во второй части - NLP, третья посвящена CV, а четвертая аудио.

Оценивание включает домашние задания в Jupyter Notebook, тесты, соревнования на платформе Kaggle и итоговый проект. Тесты автоматизированы на платформе Stepik. Домашние задания проверяются в формате кросс-валидации между студентами. Проект проверяется преподавателем, закрепленным за студентом.

Курс имеет платный тариф, в котором предлагается доступ в приватный чат, персональный ментор, приоритетная проверка домашних заданий и дополнительные вебинары.

Кросс-проверка домашних заданий является ограничением --- по тем же причинам, что и в курсе МГУ, при этом платный тариф показывает, что такой подход проверки может быть долгим. Также курс полностью не автоматизирован тестами и кросс-проверками: проект вручную проверяется преподавателем, а для более расширенной обратной связи необходимо перейти на платный тариф -- это показывает важность обратной связи и подкрепляет сложность масштабирования на большую аудиторию студентов.

\subsection{Open ML Course: Классические модели ML}

Курс сообщества Open Data Science~\cite{ods_ml_course} представляет собой полностью онлайн-формат обучения и распространяется бесплатно без отбора слушателей. Темы включают линейные модели, решающие деревья и дополнительные главы по машинному обучению. Длительность курса составляет два месяца.

Проверочные контроли состоят из тестов, дополнительно дается возможность пройти соревнование. Проверка заданий полностью автоматизирована на собственной платформе ODS.ai.

Полная автоматизация проверки достигается здесь за счёт отказа от заданий открытого
типа --- контроль сводится к тестам и соревнованиям, поэтому этот компромисс
фиксируется как нежелательный.

\subsection{Курс машинного обучения girafe-ai}

Открытый курс girafe-ai~\cite{girafe_ai} содержит полный цикл изучения машинного обучения: от линейных моделей и метода опорных векторов до нейронных сетей, алгоритма обратного распространения ошибки, эмбеддингов и архитектур seq2seq. Материалы курса и видеозаписи лекций опубликованы в открытом доступе.

Контроль знаний организован через серию заданий и лабораторных работ в Jupyter Notebook, а также итоговый экзамен. Некоторые задания сразу включают тесты для самопроверки, а некоторые содержат теоретические вопросы. Нигде явно не описано, как именно и кем проверяются работы.

Теоретические вопросы и итоговый экзамен воспроизводят
модель академических курсов и дополнительно подтверждают их подход:
открытые задания необходимы для диагностики понимания алгоритмов, а не
только их практического применения.

\subsection{Karpov.Courses: Инженер машинного обучения}

Коммерческий курс~\cite{karpov_ml} реализован полностью в онлайн-формате. Программа ориентирована на подготовку специалистов от базового уровня до внедрения моделей в производственную эксплуатацию: линейные модели, деревья решений, нейронные сети, обработка естественного языка, рекомендательные системы, A/B-тестирование и развертывание моделей.

Классические семинары в курсе отсутствуют, практическая подготовка обеспечивается за счет решения прикладных задач на собственной платформе и разработки итогового проекта в формате ML-сервиса с проверкой экспертов.

Курс демонстрирует, что проектная модель оценивания с развёрнутой обратной
связью является ключевой для подготовки специалистов. 
Полное отсутствие теоретических открытых вопросов рассматривается как
ограничение: такой подход не позволяет диагностировать глубину понимания
алгоритмов, что становится ограничением для академического курса. 

\subsection{Яндекс Практикум: ML-инженер с опытом}

Курс «ML-инженер с опытом» от Яндекс Практикума~\cite{yandex_practicum_ml} реализован полностью в онлайн-формате и рассчитан на четыре месяца обучения. Программа ориентирована на специалистов с базовыми знаниями классического машинного обучения и охватывает полный производственный цикл. 

Система оценивания строится на пяти учебных проектах и итоговом проекте, выполняемом в условиях, приближённых к реальным рабочим задачам. Проверка проектов осуществляется ревьюерами с предоставлением развёрнутой письменной обратной связи. Практика выполняется во встроенной инфраструктуре платформы с доступом к Yandex Cloud. Про теоретические задания отдельно ничего не говорится. Студентам доступен Практикум ИИ - нейросеть, отвечающая на вопросы по материалу курса.

Курс подтверждает масштабируемость проектной модели в коммерческом сегменте, как и Karpov.Courses. 

\subsection{Обобщение результатов анализа}

Результаты анализа систематизированы в трёх таблицах:
в таблице~\ref{tab:courses_general}~--- общие характеристики
и тематическое содержание программ,
в таблице~\ref{tab:courses_assessment}~--- виды оценочных мероприятий,
в таблице~\ref{tab:courses_grading}~--- организация проверки учебных заданий.

\begin{table}[h!]
\caption{Общая характеристика и тематическое содержание курсов}
\label{tab:courses_general}
\centering
\footnotesize
\setlength{\tabcolsep}{3pt}
\begin{tabular}{|p{2.5cm}|p{2.1cm}|c|c|c|c|c|p{2.2cm}|}
\hline
\textbf{Курс} &
\textbf{Тип} &
\textbf{Формат} &
\textbf{Год} &
\textbf{Кл.\,МО} &
\textbf{Гл.\,об.} &
\textbf{Внедр.} &
\textbf{Доп.\,темы} \\
\hline
МО ФКН НИУ ВШЭ   & Академический & Очный  & 2026 & Да   & ---  & ---  & --- \\
\hline
МО ВМК МГУ       & Академический & Очный  & 2026 & Да   & ---  & ---  & --- \\
\hline
Open ML Course   & Открытый      & Онлайн & 2026 & Да   & ---  & ---  & --- \\
\hline
DLS МФТИ         & Открытый      & Онлайн & 2026 & Да   & Да   & ---  & NLP, CV, аудио \\
\hline
girafe-ai        & Открытый      & Очный  & 2021 & Да   & Да   & ---  & seq2seq, эмбеддинги \\
\hline
Karpov.Courses   & Коммерческий  & Онлайн & 2026 & Да   & Да   & Да   & A/B тесты \\
\hline
Яндекс Практикум & Коммерческий  & Онлайн & 2025 & ---  & ---  & Да   & RecSys, NLP \\
\hline
\multicolumn{8}{|l|}{\textit{Кл.\,МО~--- классическое МО; Гл.\,об.~--- глубинное обучение; Внедр.~--- внедрение моделей в производство;}} \\
\multicolumn{8}{|l|}{\textit{<<--->>~--- отсутствует.}} \\
\hline
\end{tabular}
\end{table}

\begin{table}[h!]
\caption{Виды оценочных мероприятий}
\label{tab:courses_assessment}
\centering
\footnotesize
\setlength{\tabcolsep}{3pt}
\begin{tabular}{|p{2.5cm}|p{2.3cm}|c|c|c|c|c|}
\hline
\textbf{Курс} &
\textbf{Практ.\,задания} &
\textbf{Теор.\,задания} &
\textbf{Тесты} &
\textbf{Экз./КР} &
\textbf{Проект} &
\textbf{Соревн.} \\
\hline
МО ФКН НИУ ВШЭ   & Да (JN)    & Да  & Да   & Да  & ---  & ---  \\
\hline
МО ВМК МГУ       & Да (JN)    & Да  & Да   & Да  & ---  & Да   \\
\hline
Open ML Course   & ---        & --- & Да   & --- & ---  & Да   \\
\hline
DLS МФТИ         & Да (JN)    & Да  & Да   & --- & Да   & Да   \\
\hline
girafe-ai        & Да (JN)    & Да  & Да   & Да  & ---  & ---  \\
\hline
Karpov.Courses   & Да (плат.) & --- & ---  & --- & Да   & ---  \\
\hline
Яндекс Практикум & Да (плат.) & --- & ---  & --- & Да   & ---  \\
\hline
\multicolumn{7}{|l|}{\textit{JN~--- Jupyter Notebook; плат.~--- собственная платформа курса;
Соревн.~--- соревнования.}} \\
\hline
\end{tabular}
\end{table}

\begin{table}[h!]
\caption{Организация проверки домашних заданий}
\label{tab:courses_grading}
\centering
\footnotesize
\setlength{\tabcolsep}{3pt}
\begin{tabular}{|p{2.5cm}|p{3.0cm}|p{2.3cm}|c|p{2.4cm}|}
\hline
\textbf{Курс} &
\textbf{Платформа ДЗ} &
\textbf{Автотесты} &
\textbf{Кросс-пров.} &
\textbf{Ручная пров.} \\
\hline
МО ФКН НИУ ВШЭ   & Anytask          & ---    & ---  & Да \\
\hline
МО ВМК МГУ       & Собственная      & Да     & Да   & Да \\
\hline
Open ML Course   & ODS.ai           & Да     & ---  & --- \\
\hline
DLS МФТИ         & Stepik, Kaggle   & Да     & Да   & Да  \\
\hline
girafe-ai        & GitHub           & Да     & ---  & --- \\
\hline
Karpov.Courses   & Собственная      & Да     & ---  & Да  \\
\hline
Яндекс Практикум & Собственная      & ---    & ---  & Да  \\
\hline
\end{tabular}
\end{table}

Проведённый анализ позволяет сформулировать ряд ключевых наблюдений.

\textbf{Jupyter Notebook как стандарт практических заданий.}
Пять из семи рассмотренных курсов используют Jupyter Notebook
в качестве основной среды выполнения практических заданий.
В ряде программ код в ноутбуках уже проверяется автоматически
посредством юнит-тестов, что подтверждает принципиальную
осуществимость автоматизации в данной среде и обозначает Jupyter Notebook
как естественную точку интеграции для расширенной системы оценивания.

\textbf{Классическое машинное обучение и теоретические задания.}
Классические алгоритмы машинного обучения образуют обязательное ядро
всех семи программ вне зависимости от их типа и формата.
В академических курсах изучение классических алгоритмов неизменно
сопровождается теоретическими заданиями открытого типа~---
вопросами с развёрнутым ответом, интерпретацией результатов,
формулировкой выводов,~--- которые направлены на диагностику
глубины понимания математической природы алгоритмов,
а не только умения применять их на практике.
Задания открытого типа присутствуют в пяти из семи программ
и характерны прежде всего для академических и открытых курсов.

\textbf{Итоговый проект как маркер практической ориентации.}
Проект предусмотрен в трёх из семи курсов: в двух коммерческих
(Karpov.Courses, Яндекс Практикум) и одном открытом (DLS МФТИ).
Примечательно, что именно коммерческие программы последовательно
выстраивают обучение вокруг проекта как сквозной формы работы~---
это отражает запрос рынка труда на специалистов, способных решать
комплексные задачи в производственных условиях.
Вместе с тем оценивание проекта требует развёрнутой
индивидуальной обратной связи от преподавателей или ревьюеров,
что при увеличении числа студентов существенно нагружает
преподавательский состав.

\textbf{Проверка открытых заданий как ключевое узкое место.}
Оценивание теоретических заданий открытого типа сводится
к трём подходам: ручная проверка преподавателем,
кросс-проверка между студентами или исключение подобных заданий
из системы контроля.
Каждый из них имеет существенные ограничения:
ручная проверка не масштабируется при росте потока студентов,
кросс-проверка не обеспечивает согласованности оценивания
и перекладывает нагрузку на самих обучающихся,
а отказ от открытых заданий снижает глубину диагностики знаний.
Таким образом, автоматизация охватывает проверку кода
и формализованных ответов, но практически не применяется
к оценке текстовых ответов, требующих содержательного анализа.
Это противоречие определяет необходимость рассмотрения
существующих инструментов автоматической проверки.

\section{Анализ инструментов для автоматической проверки}

Задача автоматизации проверки учебных заданий решается различными инструментами: от классических систем юнит-тестирования Jupyter Notebook до коммерческих платформ и исследовательских прототипов на основе больших языковых моделей. Ниже рассмотрены наиболее распространённые решения.

\subsection{nbgrader}

nbgrader~--- открытая система для создания и проверки заданий в формате Jupyter Notebook, разработанная в UC Berkeley и поддерживаемая сообществом Jupyter~\cite{nbgrader, nbgrader_github}. Поддерживает два типа ячеек для проверки: автоматически проверяемые ячейки кода (на основе assert-тестов) и ячейки с открытым ответом. Открытые ответы (теоретические вопросы, текстовые объяснения) проверяются вручную через веб-интерфейс Formgrader. Преподавателю необходимо разметить ячейки в Jupyter, написать тесты и настроить структуру каталогов. Nbgrader доступен как в JupyterHub, так и как консольная утилита.

\subsection{Otter-Grader}

Otter-Grader~--- открытая система автоматической проверки на Python и R~\cite{otter_grader, otter_grader_github}. Как и nbgrader, использует разметку ноутбука специальными комментариями и предоставляет консольную утилиту, однако ключевое отличие состоит в модели исполнения: команда otter assign генерирует из исходного ноутбука студенческую версию и архив для запуска проверки в изолированном Docker-контейнере, что упрощает масштабирование. Открытые вопросы помечаются как manual: true и экспортируются в PDF для ручной проверки во внешней системе. Собственного интерфейса для проверки открытых ответов нет.

\subsection{CodeGrade}

CodeGrade~--- коммерческая платформа, разработанная в Университете Амстердама~\cite{codegrade}. Имеет нативную поддержку Jupyter Notebook, позволяющую визуально настраивать блоки тестов: проверки ввода-вывода, юнит-тесты, линтеры, структурные проверки  и пользовательские скрипты. Интегрирован с LMS, включает встроенную проверку на плагиат. У студентов есть доступ к ИИ-ассистенту, но пока этот функционал находится на бета-тестировании.
\subsection{PyEvalAI}

PyEvalAI~--- исследовательский прототип находящийся в апробации и без открытого кода~\cite{pyevalai}. Система оценивает как код, так и текстовые ячейки в Jupyter Notebook, комбинируя юнит-тесты с локально развёрнутой открытой БЯМ. Модель получает на вход условие задачи, эталонное решение, ответ студента и результаты тестов, после чего генерирует численную оценку и персонализированную обратную связь. От преподавателя требуется составить ноутбуки с заданиями и эталонами, а также развернуть сервер с GPU.

\subsection{Яндекс.Контест}

Яндекс.Контест~--- проприетарная платформа компании Яндекс, реализующая классическую модель онлайн-судьи (stdin/stdout)~\cite{yandex_contest}. Поддерживает программные задачи с тестами и пользовательские проверки. Jupyter Notebook не поддерживается; открытые теоретические вопросы оцениваться не могут. Преподавателю необходимо подготовить условие и тесты.

\subsection{Stepik}

Stepik~--- образовательная платформа, поддерживающая 18 типов автоматически проверяемых заданий: тесты с выбором ответа, численные задачи, проверка формул через SymPy, регулярные выражения, исполнение кода против тестов, SQL~\cite{stepik}. Для открытых ответов доступна кросс-валидация между студентами или проверка командой курса. Jupyter Notebook не поддерживается~--- код выполняется во встроенном редакторе. Преподаватель настраивает задания через веб-интерфейс, для нестандартных проверок может написать функцию check(reply) на Python.

\subsection{GitHub Classroom}

GitHub Classroom~--- бесплатный сервис GitHub для управления учебными заданиями, использующий инфраструктуру Git и GitHub Actions для автоматической проверки~\cite{github_classroom}. Каждое задание представлено шаблонным репозиторием: при принятии задания студент получает приватную копию, и при каждом \texttt{git push} в Linux-окружении GitHub Actions запускаются преподавательские тесты, результаты которых сразу отображаются в интерфейсе репозитория. .

Открытые теоретические вопросы платформа не оценивает~--- проверка строится исключительно на правилах. Jupyter Notebook нативно не поддерживается, однако ноутбук может быть преобразован в Python-скрипт или проверен через Python библиотеки.

\subsection{Kaggle InClass}

Kaggle InClass~--- режим закрытых учебных соревнований на платформе Kaggle, позволяющий преподавателю создать соревнование с приватным тестовым набором~\cite{kaggle_inclass}. Студенты загружают файлы предсказаний, которые автоматически оцениваются по выбранной метрике с публикацией результатов на лидерборде. Проверяется только качество предсказаний~--- ни код, ни теоретические ответы платформа не анализирует, что требует комбинирования с другими инструментами для полноценного контроля. Преподавателю необходимо подготовить датасет, разметку, ground truth для скрытой части и описание метрики.


\subsection{Сравнительный анализ}

В таблице~\ref{tab:grading_tools} приведено сравнение рассмотренных инструментов по ключевым для настоящей работы критериям.

\begin{table}[h!]
\LeftCaption
\caption{Сравнение инструментов автоматической проверки}
\label{tab:grading_tools}
\centering
\begin{tabular}{|c|c|c|}
\hline
\textbf{Инструмент} & \textbf{Открытые вопросы} & \textbf{Jupyter Notebook}  \\
\hline
nbgrader & Только вручную & Нативная  \\
\hline
Otter-Grader & Только вручную & Нативная \\
\hline
CodeGrade & Только вручную & Нативная \\
\hline
GitHub Classroom & Не поддерживаются & Нет  \\
\hline
PyEvalAI & Да, через БЯМ & Нативная \\
\hline
Яндекс.Контест & Не поддерживаются & Нет \\
\hline
Stepik & Только вручную или кросс-валидация & Нет \\
\hline
Kaggle InClass & Не поддерживаются & Нет \\
\hline
\end{tabular}
\end{table}

\begin{itemize}
    \item 7 из 8 инструментов либо проверяют открытые вопросы только вручную,
    либо не поддерживают их вовсе.

    \item Нативная поддержка Jupyter Notebook реализована у 4 из 8 инструментов;
    остальные 3 её не имеют.

    \item Единственный инструмент, сочетающий автоматическую проверку открытых
    вопросов через БЯМ и нативную поддержку Jupyter Notebook,~--- PyEvalAI~---
    остаётся исследовательским прототипом с ограниченной апробацией.

    \item Таким образом, 6 из 7 рассмотренных решений не закрывают потребность
    в автоматической проверке открытых теоретических вопросов в среде
    Jupyter Notebook.
\end{itemize}


\section{Итоги анализа}

Проведённый анализ существующих курсов и инструментов автоматической проверки позволяет сформулировать ключевые наблюдения.

Во-первых, задания открытого типа — теоретические вопросы с развёрнутым ответом, интерпретация результатов, формулировка выводов — присутствуют в большинстве академических курсов по машинному обучению, однако их проверка остаётся либо ручной, либо вынесенной на кросс-валидацию между студентами. Оба подхода плохо масштабируются и не обеспечивают своевременной и согласованной обратной связи.

Во-вторых, среди восьми рассмотренных инструментов автоматической проверки ни один не сочетает нативную поддержку Jupyter Notebook с автоматической оценкой открытых текстовых ответов. Единственное исключение — исследовательский прототип PyEvalAI — не имеет открытого кода и находится на стадии апробации.

В-третьих, наиболее распространённой средой для практических заданий по машинному обучению остаётся Jupyter Notebook: 5 из 6 проанализированных курсов используют именно этот формат. При этом nbgrader — наиболее зрелый инструмент с нативной поддержкой Jupyter Notebook — не предоставляет средств для автоматической оценки открытых ответов, ограничиваясь ручной проверкой через интерфейс Formgrader.

Таким образом, в существующей экосистеме инструментов для обучения машинному обучению сохраняется незакрытая потребность: автоматическая проверка открытых теоретических вопросов и проектных работ в среде Jupyter Notebook. Именно это противоречие определяет направление настоящей работы.

\section{Формулировка гипотезы и образ решения}

\subsection{Формулировка гипотезы}

На основании проведённого анализа формулируется следующая гипотеза: интеграция большой языковой модели в систему проверки заданий на базе nbgrader позволит автоматизировать оценку открытых теоретических вопросов в среде Jupyter Notebook, обеспечив при этом качество и согласованность обратной связи, сопоставимые с ручной проверкой преподавателя, и не требуя от него дополнительных временных затрат на рутинное оценивание. Помимо этого, применение БЯМ для проверки итоговых проектов позволит обеспечить студентам персонализированную развёрнутую обратную связь по нетиповым открытым заданиям.

Проверка гипотезы предполагает разработку соответствующего программного решения, создание курса по классическому машинному обучению и проведение апробации на реальной учебной аудитории.

\subsection{Образ решения}

Предлагаемое решение состоит из двух взаимосвязанных компонентов: расширения системы nbgrader с поддержкой автоматической проверки открытых ответов через БЯМ и системы проверки итогового проекта курса.

Первый компонент представляет собой расширение стандартного функционала nbgrader. В текущем виде инструмент поддерживает два типа ячеек: автоматически проверяемые ячейки с кодом на основе assert-тестов и ячейки с открытым ответом, проверяемые вручную. Предлагается ввести третий тип — ячейки с открытым ответом, оцениваемые автоматически с привлечением БЯМ. Для каждой такой ячейки преподаватель указывает условие задания, эталонный ответ и критерии оценивания. При проверке модель получает эти данные вместе с ответом студента и возвращает числовую оценку и содержательную обратную связь. Такой подход позволяет автоматизировать оценку коротких теоретических вопросов, формулировок выводов и интерпретаций результатов — тех типов заданий, которые принципиально важны для глубокого понимания материала, но традиционно требуют ручной проверки. Интеграция выполняется в рамках существующего интерфейса nbgrader, что не требует от преподавателей изменения привычного рабочего процесса.

Второй компонент — система проверки итогового проекта. В отличие от стандартных заданий проект представляет собой развёрнутую самостоятельную работу студента: постановку задачи, анализ данных, обучение и сравнение моделей, интерпретацию результатов и выводы. Проект сдаётся с кодом и текстовыми пояснениями в открытом виде. БЯМ анализирует как программную, так и текстовую части работы: корректность реализации, качество анализа, обоснованность выводов и полноту изложения. По результатам анализа модель формирует развёрнутую персонализированную обратную связь по каждому разделу работы и дает оценку. Данный подход позволяет предоставить студенту детальный разбор его работы без привлечения преподавателя для проверки каждого проекта индивидуально, а также расширить перечень решаемых задач за счёт нетиповых и открытых постановок.

Для апробации предложенных инструментов разрабатывается курс по классическому машинному обучению, охватывающий базовые разделы дисциплины: линейные модели, решающие деревья, ансамблевые методы и обучение без учителя. Структура курса предусматривает домашние задания в формате Jupyter Notebook с ячейками кода, открытыми теоретическими вопросами и итоговым проектом. Апробация проводится на студентах НИУ ВШЭ, что позволяет оценить работу системы в условиях реального учебного процесса.